% !TeX program = pdflatex
% Self-contained source. Compile twice with pdfLaTeX.
\documentclass[11pt,letterpaper]{article}
\usepackage[T1]{fontenc}
\usepackage[utf8]{inputenc}
\usepackage{lmodern,amsmath,amssymb}
\usepackage[margin=25mm]{geometry}
\usepackage{array,longtable,booktabs,calc}
\usepackage[hidelinks,unicode]{hyperref}
\providecommand{\tightlist}{\setlength{\itemsep}{0pt}\setlength{\parskip}{0pt}}
\providecommand{\pandocbounded}[1]{#1}
\setcounter{secnumdepth}{0}
\setlength{\emergencystretch}{2em}
\setlength{\parindent}{0pt}
\setlength{\parskip}{4pt plus 1pt minus 1pt}
\raggedbottom
\hypersetup{pdftitle={Swap pairs on spheres of revolution: what the equivariant spectrum sees},pdfauthor={@Erina}}
\begin{document}
{\LARGE\bfseries Swap pairs on spheres of revolution\par}
{\Large What the equivariant spectrum sees\par}
{\small @Erina\hfill November 10, 2026\par}
\medskip

Fix three nonzero nonnegative bumps with separated supports on the northern half of a round
sphere. Each bump can be placed at its given latitude or at its reflection. The equivariant
Laplace spectrum distinguishes these arrangements up to a simultaneous reflection. All
coefficients in the local semiclassical expansion agree across this family. The distinction
comes from exact inverse traces and from the different lengths assigned to a gap by meridian
arc length and by the Mercator coordinate.

The rotation-invariant sector alone also detects some triples for which the first three
inverse traces coincide. For any prescribed finite number of traces, however, a family with
enough bumps admits two distinct arrangements that agree on all the prescribed initial
traces.

\section{1. Setting and notation}\label{setting-and-notation}

Normalize the area of the sphere to \(4\pi\) and the period of the effective rotation action
to \(2\pi\). In the area coordinate \(-1<x<1\) and angle \(\theta\), a rotation-invariant
metric and the unit round profile are given by

\[g_v=v(x)\,dx^2+v(x)^{-1}\,d\theta^2,\qquad v_0(x)=(1-x^2)^{-1}.\]

Here \(v>0\) and \(v-v_0\) extends smoothly to the closed interval. These endpoint
conditions give smooth poles. Write \(f^\#(x)=f(-x)\). For non-round metrics in this
normalization, two profiles define isometric metrics precisely when they coincide or are
related by this reflection. Indeed, a nonconstant curvature function has a regular point.
The orbit there has dimension at most one. An orientation-preserving isometry fixing the
point fixes its nonzero curvature gradient, hence its derivative is the identity and the
isometry is the identity. Hence the rotation circle is the identity component of the
isometry group. Every isometry carries the latitude orbits to latitude orbits and preserves
the area between them, which gives \(x\mapsto x\) or \(x\mapsto -x\).

\textbf{Sectors.} On functions \(u(x)e^{im\theta}\) the nonnegative Laplacian acts on
\(L^2((-1,1),dx)\) by

\[L_m u=-\left(\frac{u'}{v}\right)'+m^2vu,\qquad P_v=L_0.\]

The equivariant spectrum records every eigenvalue with its angular weight \(m\) and
multiplicity. We call the labelled sectors with \(|m|>K\) the high-weight sectors. The
\(K\)-equivariant spectrum in [1, Definition 2.1] additionally retains the full ordinary
spectrum with multiplicities.

\textbf{Three meridian coordinates.} In each case choose the origin at the equator.
Distances across a gap in Sections 8--10 are computed using the indicated round coordinate.

{% Unnumbered results table.
\begin{longtable}[]{@{}
  >{\raggedright\arraybackslash}p{(\linewidth - 4\tabcolsep) * \real{0.3333}}
  >{\raggedright\arraybackslash}p{(\linewidth - 4\tabcolsep) * \real{0.3333}}
  >{\raggedright\arraybackslash}p{(\linewidth - 4\tabcolsep) * \real{0.3333}}@{}}
\toprule\noalign{}
\begin{minipage}[b]{\linewidth}\raggedright
Coordinate
\end{minipage} & \begin{minipage}[b]{\linewidth}\raggedright
Differential
\end{minipage} & \begin{minipage}[b]{\linewidth}\raggedright
On the round sphere
\end{minipage} \\
\midrule\noalign{}
\endhead
\bottomrule\noalign{}
\endlastfoot
Area & \(dx\) & \(x\) \\
Meridian arc length & \(\sqrt v\,dx\) & \(\arcsin x\) \\
Mercator & \(v\,dx\) & \(\operatorname{arctanh}x\) \\
\end{longtable}
}

\textbf{Single wells and swaps.} We consider profiles with a nondegenerate minimum at
\(x=0\) and nonvanishing derivative on either open side. For a level \(s>\min v\), write the
sublevel interval as \((\ell(s),r(s))\), its width as \(W=r-\ell\) and its midpoint as
\(a=(r+\ell)/2\). The two slope combinations used in [1, Section 3] determine \(W\) and
\(|a'|\). Choosing which inverse slope belongs to each side requires a consistent labelling
over all levels. For the bump families below certain changes of labelling amount to
reflecting one bump at fixed area coordinates.

\textbf{Bump families.} Let \(\beta_1,\ldots,\beta_N\) be smooth, nonnegative, nonzero
functions whose supports lie in ordered compact subintervals of \((0,1)\) separated by
positive gaps. Extend them by zero and choose \(\varepsilon>0\). For signs
\(\sigma_i=\pm1\), put bump \(i\) in the north for \(\sigma_i=1\) and in the south for
\(\sigma_i=-1\). Thus

\[e=v_0+\frac{\varepsilon}{2}\sum_i(\beta_i+\beta_i^\#),\qquad O_i=(\beta_i-\beta_i^\#)/2,\]

\[v_\sigma=e+\varepsilon\sum_i\sigma_i O_i,\qquad v_{-\sigma}=v_\sigma^\#.\]

For sufficiently small \(\varepsilon\), the affected height ranges are disjoint, the
branches stay strictly monotone, and the curvature \(-(1/v_\sigma)''/2\) stays positive. At
each height the two endpoints of the sublevel interval for any two arrangements agree up to
reflection. Thus the arrangements share \(e\), \(W\) and \(|a'|\). We ask whether spectral
data distinguish the sign vectors up to their common reversal.

\section{2. A smooth swap and the inverse branch
argument}\label{a-smooth-swap-and-the-inverse-branch-argument}

Dryden, Macedo and Sena-Dias [1, Theorem 1.4] state that the asymptotic equivariant
spectrum determines the metric in their single-well class. We examine the last step of the
proof in the authors' arXiv v2 [1, Section 3, p.~16]. Fix two bumps \(A,B\) as in
Section 1, with \(A\) closer to the equator, and set

\[v_1=v_0+\varepsilon(A+B),\qquad v_2=v_0+\varepsilon(A+B^\#).\]

\textbf{Proposition 2.1.} For all sufficiently small \(\varepsilon>0\), these profiles
define non-isometric, positively curved single-well metrics. They agree with the round
metric near the equator and both poles. Every coefficient in the expansion of the smoothed
equivariant trace in [1, Theorem 2.2] takes the same value on both profiles. In
particular, their invariants (6) and (7) of [1, Theorem 2.5] agree for every test
function and admissible nonzero weight parameter in that theorem.

\textbf{Proof.} On the union of the northern support intervals \(v_0'\) has a positive lower
bound. Small perturbations preserve strict monotonicity there and on the reflected
intervals. Positivity of curvature follows from \(C^2\) continuity on these compact
intervals. The origin retains its round, nondegenerate minimum. The two arrangements differ
and are not global reflections, so Section 1 gives non-isometry.

Take an interval \(I\) containing the support of \(B\) and disjoint from the support of
\(A\). Outside \(I\cup(-I)\) the profiles coincide. On \(I\cup(-I)\) one is the reflection
of the other. The densities in (6) and (7) depend locally on \(v\), \((v')^2\) and \(v''\)
and are invariant under simultaneous reflection of the position and momentum variables. More
generally, the recursion for the coefficients \(b_k\) in [1, Theorem 2.2] is local in
\(v\), its inverse and their derivatives. Each first-order differentiation changes sign
under \((x,\xi)\mapsto(-x,-\xi)\) along with its momentum factor, and the two derivative
factors in the second-order term change sign together. Starting with \(b_0=1\), the
recursion gives \(b_k[v^\#](x,\xi,t)=b_k[v](-x,-\xi,t)\). Changing variables on the
reflected intervals leaves each integrated coefficient unchanged.

\textbf{The labelling step.} Use the increasing inverse side lengths \(f_1=r\) and
\(f_2=-\ell\), so \(p=f_1'\) and \(q=f_2'\) are both positive. The quantities \(p+q\) and
\(1/p+1/q\) determine \(\{p(s),q(s)\}\) at each height. To integrate them as globally
labelled functions one must still decide which branch to follow across a height where their
jets agree. Infinite-order agreement permits a smooth change of labelling. A candidate must
also satisfy the endpoint conditions after integration. The profiles above satisfy these
conditions and realize a change of labelling which is not a global reflection.

The generic symmetry-defect hypothesis in Colin de Verdière's single-well theorem [2,
Theorem 5.1(3)] gives a model for a sufficient assumption excluding such local changes.
Real analyticity likewise gives unique continuation of the inverse branches at regular
levels. The smooth counterexample concerns recovery from the coefficients of the trace
expansion. The exact data in [1, Definition 2.1] separate these two metrics, as Sections
3, 6 and 9 demonstrate. The earlier mirror-symmetric case is treated by Brüning and Heintze
[4].

\section{3. The rotation invariant sector and its second inverse
trace}\label{the-rotation-invariant-sector-and-its-second-inverse-trace}

We omit the constant mode when taking the inverse of \(P_v\). For a mean-zero function
\(f\), define \((Sf)(x)=\int_{-1}^{x}f(t)\,dt\). Solving \(P_vu=f\) gives \(u'=-vSf\), with
the additive constant fixed by the mean-zero condition. Thus \(P_v^{-1}=S^*M_vS\) as a
quadratic form on the mean-zero subspace, where \(M_v\) denotes multiplication by \(v\).
Since \(Sf\) vanishes at both endpoints, \(SS^*\) has kernel equal to the Dirichlet Green
function for the ordinary interval Laplacian,

\[C(x,y)=\frac{(1+\min(x,y))(1-\max(x,y))}{2}.\]

With \(x_{k+1}=x_1\), cyclically composing the kernels gives

\[T_k(v)=\sum_{j\ge1}\lambda_{0,j}^{-k}=\int_{(-1,1)^k}\prod_{i=1}^k v(x_i)C(x_i,x_{i+1})\,dx_1\cdots dx_k.\]

Indeed \(M_{\sqrt v}S\) is Hilbert-Schmidt since \(\int v(x)C(x,x)\,dx\) is finite. Each
cyclic integral converges by \(C(x,y)^2\le C(x,x)C(y,y)\). In particular,

\[T_1(v)=\frac12\int_{-1}^{1}(1-x^2)v(x)\,dx,\qquad J(v)=T_2(v)=\iint v(x)v(y)C(x,y)^2\,dx\,dy.\]

The first formula also appears in [3, Proposition 3.1]. On the round sphere
\(J(v_0)=\pi^2/3-3\), as follows by summing \(1/(\ell^2(\ell+1)^2)\) over \(\ell\ge1\). For
a positive eigenvalue \(\lambda\), the flux \(w=u'/v\) satisfies the fixed-end string
equation \(-w''=\lambda v w\) with \(w(\pm1)=0\).

\textbf{Proposition 3.1.} Let \(v=e+o\) be its decomposition into even and odd parts. On the
northern half put

\[K(x,y)=C(x,y)^2-C(x,-y)^2=\min(x,y)(1-\max(x,y))^2.\]

Then \(J(v)=J(e)+2\langle o,Ko\rangle\), with the pairing over \((0,1)\). The excess is
nonnegative and vanishes precisely for an even profile.

\textbf{Proof.} The mixed even-odd integral vanishes by reflection. Folding the odd-odd
integral onto \((0,1)^2\) gives the asserted kernel. If \(G_D\) denotes the Dirichlet Green
operator on \((0,1)\), then \(K(x,y)=G_D(x,y)(1-\max(x,y))\). Write
\(o_t=o\,\mathbf 1_{(0,t)}\). Fubini gives
\(\langle o,Ko\rangle=\int_0^1\langle o_t,G_D o_t\rangle\,dt\). The inverse Dirichlet
Laplacian is positive and injective. For \(o\ne0\), the integrand is strictly positive for
all \(t\) sufficiently close to one.

\textbf{Corollary 3.2.} The two-bump profiles of Section 2 have different \(m=0\) spectra.
On the northern half their odd parts are \(\varepsilon(A+B)/2\) and \(\varepsilon(A-B)/2\).
Thus

\[J(v_1)-J(v_2)=2\varepsilon^2\langle A,KB\rangle=2\varepsilon^2\int_0^1\int_0^1 A(x)K(x,y)B(y)\,dy\,dx>0.\]

The inequality is strict because \(K\) is pointwise positive and each bump is nonzero and
nonnegative. This is an exact identity at every admissible amplitude. Their first inverse
traces agree because the weight in \(T_1\) is even. More generally, on a fixed bump family
the algebraic dependence on the signs is

\[J(v_\sigma)=J(e)+\frac{\varepsilon^2}{2}\sum_{i,j}\sigma_i\sigma_j\langle\beta_i,K\beta_j\rangle.\]

\section{4. The trace in height coordinates}\label{the-trace-in-height-coordinates}

Write \(c=\min v\). For \(s>c\) let \(I_s=(\ell(s),r(s))\), \(W=r-\ell\), \(a=(r+\ell)/2\),
and \(u=1-W/2\). The profiles considered here have nonzero derivative on either side of the
minimum, so the inverse branches are smooth for \(s>c\). In formulas below, functions
indexed by ``in'' are evaluated at the inner level \(s\) and functions indexed by ``out'' at
the outer level \(t>s\).

\textbf{Proposition 4.1.} Define the inner moments by

\[E_{\mathrm{in}}=\int_{\ell}^{r}\frac{1+x^2}{2}\,dx=\frac{W}{2}(1+a^2)+\frac{W^3}{24},\qquad O_{\mathrm{in}}=aW.\]

The corresponding outer moments are

\[E_{\mathrm{out}}=\frac{(1-r)^3+(1+\ell)^3}{6}=\frac{u^3}{3}+ua^2,\]

\[O_{\mathrm{out}}=\frac{(1+\ell)^3-(1-r)^3}{6}=au^2+\frac{a^3}{3}.\]

For \(c<s<t\), the density of the sum over the four boundary pairs is

\[\sum_{x\in\partial I_s,\ y\in\partial I_t} C(x,y)^2\,|x'(s)y'(t)|=-E_{\mathrm{in}}'(s)E_{\mathrm{out}}'(t)+O_{\mathrm{in}}'(s)O_{\mathrm{out}}'(t).\]

\textbf{Proof.} Put \(F(s,t)=\int_{I_s}\int_{I_t^c}C(x,y)^2\,dy\,dx\), with complements
taken in \((-1,1)\). On the left tail \(y<\ell(t)\) the squared kernel factors as
\((1+y)^2(1-x)^2/4\). On the right tail \(y>r(t)\) it factors as \((1+x)^2(1-y)^2/4\). The
inner integrals of \((1\pm x)^2\) are \(2(E_{\mathrm{in}}\pm O_{\mathrm{in}})\). The two
outer tail integrals are \((1+\ell(t))^3/3\) on the left and \((1-r(t))^3/3\) on the right.
Multiplication gives

\[F(s,t)=E_{\mathrm{in}}(s)E_{\mathrm{out}}(t)-O_{\mathrm{in}}(s)O_{\mathrm{out}}(t).\]

The sum in the proposition is \(-\partial_s\partial_tF\), since \(I_s\) expands with \(s\)
and \(I_t^c\) contracts with \(t\).

One can recover \(J\) by integrating the boundary-pair density against \(2st\,ds\,dt\) over
\(c<s<t\). The layer-cake formula also identifies the dependence on the signs of the
midpoints directly. Put \(I_s=\varnothing\) for \(s<c\) and set

\[D(t)=\int_{I_t^c}\frac{(1-y^2)^2}{6}\,dy.\]

Indeed, \(\int_{-1}^1 C(x,y)^2\,dx=(1-y^2)^2/6\). Thus the layer integrand over
\(I_s^c\times I_t^c\) for \(s<t\) is \(D(t)-F(s,t)\), with \(F=0\) when \(s<c\).
Consequently \(J=J_{\mathrm{sq}}+J_{\mathrm{sign}}\), where

\[J_{\mathrm{sq}}=2\int_0^\infty tD(t)\,dt-2\iint_{c<s<t}E_{\mathrm{in}}(s)E_{\mathrm{out}}(t)\,ds\,dt,\]

\[J_{\mathrm{sign}}=\frac{2}{3}\iint_{c<s<t}W(s)a(s)a(t)\bigl[3u(t)^2+a(t)^2\bigr]\,ds\,dt.\]

At fixed \(W\), the first expression depends on the path only through \(a^2\), whereas the
second couples the signs at different levels. The first inverse trace has the simpler
expression

\[T_1(v)=\int_0^\infty\left(\frac23-\frac W2+\frac{W^3}{24}+\frac{Wa^2}{2}\right)ds,\]

where \(W=a=0\) below \(c\). This follows by integrating \((1-x^2)/2\) over the complement
of each sublevel interval.

\section{5. Three bumps with equal first three inverse
traces}\label{three-bumps-with-equal-first-three-inverse-traces}

Take nonzero, nonnegative smooth bumps \(A,B,D\) with disjoint, ordered supporting intervals
in \((0,1)\). Keep \(B\) in the north, put \(rD\) in the south, and compare the two
placements of \(A\):

\[v_N=v_0+\varepsilon(A+B+rD^\#),\qquad v_S=v_0+\varepsilon(A^\#+B+rD^\#).\]

For small \(\varepsilon>0\) the profiles share \(W,|a'|\), their even part \(e\), and
\(T_1\). The fixed \(B\) prevents a global reflection. On the northern half the odd parts
are \(\varepsilon(\pm A+B-rD)/2\). By Proposition 3.1,

\[T_2(v_N)-T_2(v_S)=2\varepsilon^2\langle A,K(B-rD)\rangle.\]

Thus the unique positive balancing amplitude is
\(r=\langle A,KB\rangle/\langle A,KD\rangle\).

For the third trace, put

\[G_e(x,y)=\int_{-1}^1 C(x,z)e(z)C(z,y)\,dz,\]

\[Q_e(x,y)=C(x,y)G_e(x,y)-C(x,-y)G_e(x,-y).\]

In the triple integral for \(T_3(e+o)\), simultaneous reflection removes the terms with an
odd number of odd factors. The three terms with two odd factors coincide, and folding them
onto the northern half gives \(T_3(e+o)=T_3(e)+6\langle o,Q_eo\rangle\). Hence

\[T_3(v_N)-T_3(v_S)=6\varepsilon^2\langle A,Q_e(B-rD)\rangle.\]

\textbf{Lemma 5.1.} For \(0<x<y<1\), set \(\mathcal R_e(x,y)=Q_e(x,y)/K(x,y)\). Define

\[I_1=(1-x)\int_0^x(1+z^2/x)e(z)\,dz,\qquad I_2=\int_x^y(1-z^2)e(z)\,dz,\]

\[I_3=\frac{1+y}{1-y}\int_y^1(1-z)^2e(z)\,dz.\]

Then \(2\mathcal R_e=I_1+I_2+I_3\), and

\[\partial_y\mathcal R_e=\frac{1}{(1-y)^2}\int_y^1(1-z)^2e(z)\,dz>0,\]

\[\partial_x\mathcal R_e=-\frac12\int_0^x(1+z^2/x^2)e(z)\,dz<0.\]

\textbf{Proof.} For \(x,z>0\) the sum \(C(x,z)+C(x,-z)\) is \(1-\max(x,z)\), and their
difference is \(\min(x,z)(1-\max(x,z))\). Fold the integral defining \(G_e\) using these two
kernels and the evenness of \(e\). Splitting at \(z=x\) and \(z=y\) and dividing \(Q_e\) by
\(K=x(1-y)^2\) gives the three displayed integrals. Differentiation cancels the endpoint
terms.

\textbf{Proposition 5.2.} Choose \(r\) to balance \(T_2\). The sign of
\(\Delta T_3=T_3(v_N)-T_3(v_S)\) is fixed in the two exterior orderings. The table also
gives the interior construction of Proposition 5.3.

{% Unnumbered results table.
\begin{longtable}[]{@{}
  >{\raggedright\arraybackslash}p{(\linewidth - 4\tabcolsep) * \real{0.3333}}
  >{\raggedright\arraybackslash}p{(\linewidth - 4\tabcolsep) * \real{0.3333}}
  >{\raggedright\arraybackslash}p{(\linewidth - 4\tabcolsep) * \real{0.3333}}@{}}
\toprule\noalign{}
\begin{minipage}[b]{\linewidth}\raggedright
Support order
\end{minipage} & \begin{minipage}[b]{\linewidth}\raggedright
Third-trace difference
\end{minipage} & \begin{minipage}[b]{\linewidth}\raggedright
Comparison
\end{minipage} \\
\midrule\noalign{}
\endhead
\bottomrule\noalign{}
\endlastfoot
\(A<B<D\) & \(\Delta T_3<0\) & Increase the larger variable. \\
\(B<D<A\) & \(\Delta T_3>0\) & Increase the smaller variable. \\
\(B<A<D\) & Both signs occur in the construction. & Move the middle bump. \\
\end{longtable}
}

\textbf{Proof.} Compare the two \(K\)-weighted averages of \(\mathcal R_e\) for the pairs
\((A,B)\) and \((A,D)\). For \(A<B<D\), the factorization \(K(x,y)=x(1-y)^2\) gives the same
normalized \(A\)-weight in both averages. Holding the \(A\) variable fixed, moving the
second variable from \(B\) to \(D\) strictly increases \(\mathcal R_e\). For \(B<D<A\), the
common \(A\)-weight is proportional to \((1-y)^2 A(y)\), and moving the smaller argument
from \(B\) to \(D\) strictly decreases \(\mathcal R_e\).

\textbf{Proposition 5.3.} In the order \(B<A<D\) one can choose smooth bumps so that
\(T_1,T_2,T_3\) all agree. They can be chosen with the gap from \(A\) to \(D\) strictly
shorter than the gap from \(B\) to \(A\) in the round latitude coordinate.

\textbf{Proof.} First freeze the even profile at \(v_0\) and consider point concentrations
at centers \(b<a<d\). The difference of the two normalized averages is then

\[H(a)=\mathcal R_{v_0}(b,a)-\mathcal R_{v_0}(a,d).\]

For this limiting comparison \(H\) is strictly increasing, negative at \(b\) and positive at
\(d\), by Lemma 5.1 and continuity on the diagonal. We locate a sign change with a strict
margin in the gap inequality. Choose \(b=\delta\), \(d=3\delta\), and \(a_-=5\delta/2\). On
\(b\le x\le y\le d\), the formulas in the lemma give
\(\partial_x\mathcal R_{v_0}=O(\delta)\) and
\(\partial_y\mathcal R_{v_0}=2\log 2-1+O(\delta)\) uniformly. Subtracting through
\(\mathcal R_{v_0}(a_-,a_-)\) therefore gives

\[H(a_-)=-\tfrac{\delta}{2}(2\log 2-1)+O(\delta^2)<0.\]

Fix \(\delta\) small, and choose \(a_+\in(a_-,d)\) close enough to \(d\) that \(H(a_+)>0\).
Throughout \([a_-,a_+]\), the \(x\)-gap from \(a\) to \(b\) is at least \(3\delta/2\),
whereas the gap from \(a\) to \(d\) is at most \(\delta/2\). The round latitude density
\((1-x^2)^{-1/2}\) and Mercator density \((1-x^2)^{-1}\) are both \(1+O(\delta^2)\) on this
interval and the adjoining gaps, so these inequalities remain strict in both coordinates.

Replace the point concentrations by nonnegative smooth bumps of unit mass with sufficiently
narrow supports, translating the middle bump with \(a\). At \(\varepsilon=0\) the two
normalized averages still have opposite difference signs at \(a_-,a_+\), by continuity of
the kernels near the separated point pairs. The support gaps remain disjoint and satisfy the
strict latitude comparison uniformly on \([a_-,a_+]\). For these fixed widths the balancing
\(r(a)\) is positive, continuous and bounded. Now choose \(\varepsilon>0\) small enough that
the endpoint signs persist and all profiles have strictly positive curvature. The
intermediate value theorem applied to the actual normalized difference with its varying
\(e=e(a,\varepsilon)\) gives a position with equal \(T_3\). The choice of \(r(a)\) keeps
\(T_2\) equal throughout.

\section{6. Ordinary multiplicities and finite trace
coincidences}\label{ordinary-multiplicities-and-finite-trace-coincidences}

\textbf{Proposition 6.1 (Engman [3, Proposition 3.1]).} The rotation-invariant spectrum
is the set of ordinary Laplace eigenvalues with odd multiplicity.

\textbf{Proof.} Each separated radial problem has simple eigenvalues. The modes \(m\) and
\(-m\) for \(m\ne0\) contribute in pairs. The contribution at \(m=0\) is either zero or one,
and determines the parity of the total multiplicity.

In particular the ordinary spectrum determines every \(T_k\). It distinguishes the pair of
Corollary 3.2 and the two orderings covered by Proposition 5.2.

\textbf{Theorem 6.2.} For every \(n\ge2\), there are nonisometric smooth, positively curved
single-well profiles with the same \((W,|a'|)\) and identical \(T_1,\ldots,T_n\). At most
\(n+1\) nonnegative bumps suffice.

\textbf{Proof.} Choose an anchor \(A\) and \(n\) further nonzero bumps \(B_i\) supported in
disjoint closed subintervals of \((0,1)\). For \(u\in S^{n-1}\) write \(u_i^+=\max(u_i,0)\)
and \(u_i^-=\max(-u_i,0)\). Set

\[v_u=v_0+\varepsilon A+\varepsilon\sum_{i=1}^n(u_i^+B_i+u_i^-B_i^\#).\]

Take \(\varepsilon\) uniformly small. The affected height ranges stay disjoint because
\(v_0\) is strictly increasing on the northern support intervals. All profiles remain single
wells with strictly positive curvature by uniform \(C^2\) control. In each affected range,
the sublevel intervals for \(u\) and \(-u\) are identical or reflected. Their \((W,|a'|)\)
agree, as do their local semiclassical coefficients. Their first inverse traces agree
because \(T_1\) has an even weight. The fixed anchor excludes global reflection, and some
\(u_i\) is nonzero, so the profiles are distinct and hence nonisometric by Section 1.

The functions \(T_j(v_u)\) are continuous in \(u\), for example by the cyclic integral
formula of Section 3. The continuous odd map \(u\mapsto(T_j(v_u)-T_j(v_{-u}))_{j=2}^n\)
takes \(S^{n-1}\) into \(\mathbb R^{n-1}\). By the Borsuk--Ulam theorem it has a zero,
giving the required pair.

The construction permits both the number of bumps and the resulting pair to change with
\(n\).

\section{7. Finite determinacy for fixed bump
shapes}\label{finite-determinacy-for-fixed-bump-shapes}

\textbf{Theorem 7.1.} Fix \(N\) bump shapes and their support positions. There is an integer
\(K_*\) such that, for all admissible nonnegative amplitudes and all choices of hemisphere,
equality of \(T_1,\ldots,T_{K_*}\) is equivalent to equality of the complete \(m=0\) spectra
of the two arrangements with those amplitudes.

\textbf{Proof.} Fix two sign patterns \(\sigma,\tau\) and denote the common amplitude vector
by \(b\). The cyclic integral for \(T_j\) is homogeneous of degree \(j\) in \(v\), and the
two profiles are affine in \(b\). Hence \(p_j(b)=T_j(v_{\sigma,b})-T_j(v_{\tau,b})\) is a
polynomial of degree at most \(j\). By the Hilbert basis theorem the ascending chain of
ideals \((p_1,\ldots,p_k)\) in \(\mathbb R[b_1,\ldots,b_N]\) stabilizes. Beyond that index
every \(p_j\) lies in the stabilized ideal, so their common zero sets coincide. Take the
maximum over the finitely many pairs of patterns. Finally the power sums of a positive
trace-class inverse determine its nonzero eigenvalues with multiplicities, by successively
extracting the largest eigenvalue and subtracting its contribution.

The proof supplies no explicit cutoff. Theorem 6.2 allows the shapes to change with \(n\).
If nontrivial configurations in one fixed finite family could be matched through every
finite cutoff, even using different amplitudes at different cutoffs, applying Theorem 7.1 at
that family's \(K_*\) would already give a fully \(m=0\)-isospectral pair in it.

There is a local analytic analogue for the nonzero-sector traces \(J_m\) of Section 9. For
fixed shapes their differences between two sign patterns extend to entire functions of the
complex amplitudes by the exponential integral formula in that section. Near any point of a
compact amplitude set, Noetherianity of the ring of convergent power series gives a finite
generating list of germs from these differences. Choose local representatives of the
finitely many irreducible components of their common complex zero germ. Each remaining
entire difference vanishes as a germ on every component, hence on each representative by
analytic continuation. Thus the same finite list describes the common zero set in a
neighborhood. A finite cover of the compact set gives a single finite list of sectors there.

\section{8. Latitude gaps and the rotation-invariant
spectrum}\label{latitude-gaps-and-the-rotation-invariant-spectrum}

The spectral statistic used in this section is defined on the mean-zero sector. For
\(\kappa>0\), let

\[L_\kappa(v)=\sum_{j\ge1}\log(1+\kappa^2/\lambda_{0,j}(v))=\log\det(I+\kappa^2P_v^{-1}).\]

Its Taylor coefficients at \(\kappa=0\) give the inverse traces \(T_j\). Write
\(G_{\kappa,v}\) for the Dirichlet Green kernel of \(-d^2/dx^2+\kappa^2v\) on \((-1,1)\).

\textbf{Lemma 8.1 (mixed variation).} Let \(e\) be even and let \(P,Q\) be smooth, compactly
supported odd perturbations. Suppose \(v_{s,t}=e+sP+tQ\) is positive and has the prescribed
pole behavior throughout \([-1,1]^2\). Write \(G=G_{\kappa,v_{s,t}}\). Then

\[\partial_s\partial_t L_\kappa(v_{s,t})=-\kappa^4\iint P(x)Q(y)G(x,y)^2\,dx\,dy.\]

\textbf{Proof.} With the integration operator \(S\) of Section 3, the resolvent identity
gives \(G=S(I+\kappa^2S^*M_{v_{s,t}}S)^{-1}S^*\) as an operator. Differentiating the
trace-class determinant gives \(\partial_tL_\kappa=\kappa^2\int Q(y)G(y,y)\,dy\).
Differentiating the resolvent gives \(\partial_s G=-\kappa^2 G M_P G\), proving the formula.
Reflection identifies opposite corners of the parameter square. Hence

\[2\big[L_\kappa(e+P+Q)-L_\kappa(e+P-Q)\big]=\int_{-1}^1\int_{-1}^1\partial_s\partial_t L_\kappa(v_{s,t})\,ds\,dt.\]

\textbf{Theorem 8.2 (unequal latitude gaps).} For the profiles \(v_N,v_S\) of Section 5,
assume \(B<A<D\) and that all three bump amplitudes are positive. Use \(\xi=\arcsin x\) to
measure the two gaps between the actual facing support edges, and denote their lengths by
\(d_{BA}\) and \(d_{AD}\). If \(d_{BA}\ne d_{AD}\), then

\[\lim_{\kappa\to\infty}-\frac{1}{2\kappa}\log\big|L_\kappa(v_N)-L_\kappa(v_S)\big|=\min(d_{BA},d_{AD}).\]

Consequently the two rotation-invariant spectra differ. Their ordinary Laplace spectra also
differ by Proposition 6.1.

\textbf{Proof.} In Lemma 8.1 use the common even part \(e\), take \(P\) to be the odd part
of \(\varepsilon(B+rD^\#)\) and \(Q\) the odd part of \(\varepsilon A\). Throughout the
parameter square, \(v_{s,t}\ge v_0\) because the bump contributions are nonnegative convex
combinations of the two hemisphere placements. On each open gap \(v_{s,t}=v_0\).

We first prove a uniform estimate for one pair of adjacent bumps \(F,H\) in the same
hemisphere. Let \(d\) be the round latitude gap between their facing support edges. With
both bump functions taken positive in sign, set

\[I_\kappa(s,t)=\iint F(x)H(y)G_{\kappa,v_{s,t}}(x,y)^2\,dx\,dy.\]

The change of variable \(\xi=\arcsin x\) with \(h(x)=v_0(x)^{1/4}\) gives
\(G(x,y)=\widehat G(\xi,\eta)/(h(x)h(y))\), where \(\eta=\arcsin y\) and \(\widehat G\) is
the Green kernel on \((-\pi/2,\pi/2)\) of

\[-\partial_\xi^2+\frac{3}{4\cos^2\xi}-\frac14+\kappa^2\frac{v_{s,t}}{v_0}.\]

The potential is at least \(\kappa^2\). Green kernel comparison with the free kernel on the
line gives \(\widehat G(\xi,\eta)\le e^{-\kappa|\xi-\eta|}/(2\kappa)\). Consequently
\(I_\kappa(s,t)\le C\kappa^{-2}e^{-2\kappa d}\) uniformly on the parameter square. The same
upper bound with the appropriate separation applies to any two supports.

For the reverse bound, fix \(\delta>0\). Enlarge the gap slightly into both supports to form
a compact interval of latitudes \(J\). The bump functions vanish at the facing edges. Thus
\(J\) can be chosen with length at most \(d+\delta\) and \(v_{s,t}/v_0\le1+\delta\) on \(J\)
for all \((s,t)\). The remaining potential on \(J\) has a finite upper bound \(C_J\). Choose
closed subintervals in the interior of \(J\) on which \(F\) and \(H\), respectively, have
positive lower bounds. Such choices are possible arbitrarily near the facing edges by the
definition of support. Domain restriction and potential comparison bound \(\widehat G\)
below on this subrectangle by the Dirichlet kernel on \(J\) for \(-\partial_\xi^2+\mu^2\),
where \(\mu^2=(1+\delta)\kappa^2+C_J\). For \(J=[l,r]\) and \(\xi<\eta\) this kernel is

\[\frac{\sinh(\mu(\xi-l))\sinh(\mu(r-\eta))}{\mu\sinh(\mu(r-l))}.\]

Since both chosen subintervals stay inside \(J\), this expression is at least
\(c_\delta\mu^{-1}e^{-\mu(d+\delta)}\) for large \(\kappa\). On this compact set the
Jacobians and \(1/h\) are bounded below by positive constants. Integrating the square gives
\(I_\kappa(s,t)\ge c'_\delta\mu^{-2}e^{-2\mu(d+\delta)}\). Taking logarithms and then
letting \(\delta\downarrow0\) gives rate \(d\), uniformly in \((s,t)\). These bounds also
hold after integration over the parameter square, with adjusted constants.

In the northern hemisphere the factors in the mixed variation are \(P=\varepsilon(B-rD)/2\)
and \(Q=\varepsilon A/2\). The \(B,A\) and \(D,A\) products have opposite signs. In the
southern hemisphere both signs reverse, preserving each product's sign. Each cross-equator
separation is strictly greater than the same pair's gap in a single hemisphere because all
the original supports are compactly contained in \((0,1)\). The two same-hemisphere
contributions from the smaller gap therefore have the same sign and a strictly smaller
exponent than every potentially cancelling contribution. Lemma 8.1 now gives the stated
limit.

\textbf{Corollary 8.3.} There are three-bump pairs with the same \((W,|a'|)\) and the same
first three rotation-invariant inverse traces, but different rotation-invariant spectra. In
the near-equator construction of Section 5, choose two bracketing positions for \(A\)
strictly on the \(D\) side of the round-latitude midpoint. With narrow supports,
\(d_{AD}<d_{BA}\) throughout that bracket, including at a zero of the third-trace
difference. Theorem 8.2 applies there. By Theorem 7.1 some \(T_k\) with \(4\le k\le K_*\)
separates that pair.

\section{9. Nonzero sectors and the Mercator
kernel}\label{nonzero-sectors-and-the-mercator-kernel}

The two-bump pair of Section 2 can be distinguished by the spectrum in each individual
sector \(m\ne0\). Throughout this section write \(M=|m|\).

\textbf{Lemma 9.1 (Green kernel).} Relative to area measure \(dy\), the kernel of
\(L_m^{-1}\) is

\[K_{v,m}(x,y)=\frac{1}{2M}\exp\left(-M\left|\int_x^y v(z)\,dz\right|\right).\]

\textbf{Proof.} In the Mercator coordinate \(\tau\), defined by \(d\tau=v\,dx\), the
operator is \(v(-\partial_\tau^2+M^2)\), and the two poles are \(\tau=\pm\infty\). The
decaying full-line Green function gives the stated formula after conversion back to \(dy\).
Its weighted flux \(v^{-1}\partial_x K_{v,m}\) has jump \(-1\) at \(x=y\), and its behavior
at each pole is the regular solution in that sector.

The first inverse trace in each nonzero sector is universal [3, Proposition 3.1]. The
second inverse trace, denoted by \(J_m\), records integrals of \(v\) over intervals:

\[\mathrm{Tr}(L_m^{-1})=\int_{-1}^{1}K_{v,m}(x,x)\,dx=1/M,\]

\[J_m(v)=\mathrm{Tr}(L_m^{-2})=\frac{1}{2M^2}\iint_{-1<x<y<1}\exp\left(-2M\int_x^y v(z)\,dz\right)\,dx\,dy.\]

On the round sphere the first formula is the telescoping sum of \(1/(\ell(\ell+1))\) over
\(\ell\ge M\).

\textbf{Theorem 9.2 (two-bump identity).} Let \(\alpha,\beta\) be two bumps as in Section 1,
with \(b=\sup\,\operatorname{supp}(\alpha)<c=\inf\,\operatorname{supp}(\beta)\). Let
\(v_1=v_0+\varepsilon(\alpha+\beta)\) and \(v_2=v_0+\varepsilon(\alpha+\beta^\#)\), and
write \(v_{1,2}=e+\varepsilon A\pm\varepsilon B\), where \(e\) is their common even part,
\(A=(\alpha-\alpha^\#)/2\) and \(B=(\beta-\beta^\#)/2\). On \(-b<x<b\) and \(c<y<1\) put

\[U(x)=\frac{\varepsilon}{2}\int_{|x|}^{b}\alpha(z)\,dz,\qquad V(y)=\frac{\varepsilon}{2}\int_c^y\beta(z)\,dz.\]

Then every nonzero sector distinguishes the two profiles through the exact identity

\[J_m(v_1)-J_m(v_2)=\frac{2}{M^2}\int_{-b}^{b}\int_c^1 e^{-2M\int_x^y e(z)\,dz}\sinh(2MU(x))\sinh(2MV(y))\,dy\,dx>0.\]

\textbf{Proof.} Pair each interval \([x,y]\) with \([-y,-x]\). Their even integrals coincide
and their odd integrals have opposite signs. Using
\(\cosh(u+w)-\cosh(u-w)=2\sinh u\,\sinh w\) leaves a product of two hyperbolic sines. The
two odd integrals are simultaneously nonzero only when the two absolute endpoint values span
both bump ranges. Each surviving reflection pair has a representative in
\((-b,b)\times(c,1)\), on which the two odd integrals, including \(\varepsilon\), are
\(U(x),V(y)\). Both are nonnegative and both are positive on a set of positive area.

\textbf{Corollary 9.3 (rate).} Set \(\Delta=\int_b^c v_0(z)\,dz\), the round Mercator gap
between the true facing support edges, and let \(\delta J_m=J_m(v_1)-J_m(v_2)\). For every
fixed admissible \(\varepsilon>0\),

\[0<\delta J_m\le C M^{-2}e^{-2M\Delta},\qquad \lim_{M\to\infty}-\frac{\log\delta J_m}{2M}=\Delta.\]

\textbf{Proof.} The right-hand side of the identity in Theorem 9.2 can be rewritten as

\[\frac{1}{2M^2}\int_{-b}^b\int_c^1 e^{-2M\int_x^y v_1^\#(z)\,dz}(1-e^{-4MU(x)})(1-e^{-4MV(y)})\,dy\,dx.\]

The density \(v_1^\#\) is positive and equals \(v_0\) on \([b,c]\), which gives the upper
bound. Given \(\eta>0\), choose a small closed rectangle with \(x\) just below \(b\) and
\(y\) just above \(c\). The rectangle can be chosen so that \(\int_x^y v_1^\#<\Delta+\eta\)
while \(U\) and \(V\) have strictly positive lower bounds. Its contribution is at least
\(c_\eta M^{-2}e^{-2M(\Delta+\eta)}\) for all sufficiently large \(M\). Sending \(\eta\) to
zero gives the limit. This is a logarithmic rate for one exact statistic. Equality of all
coefficients in the smoothed trace expansion was proved separately in Proposition 2.1.

\textbf{Autocorrelation.} Let \(\psi(\tau)=1/v(x(\tau))\), so that \(dx=\psi(\tau)\,d\tau\),
and let \(R_\psi(h)=\int_{\mathbb R}\psi(t)\psi(t+h)\,dt\). This continuous even function is
integrable. The trace formula becomes

\[4M^2J_m(v)=\int_{\mathbb R}e^{-2M|h|}R_\psi(h)\,dh.\]

The high-weight values determine \(R_\psi\). Indeed, push the finite measure
\(2R_\psi(h)\,dh\) on \(h\ge0\) forward by \(q=e^{-2h}\) to \([0,1]\). The right-hand side
gives all sufficiently high integer moments of this measure. Multiplication by \(q^{M_0}\),
with \(M_0\) above the cutoff, yields a measure with every moment known. Polynomials are
dense in \(C[0,1]\) and the original pushforward has no atom at \(q=0\), so it is uniquely
determined. Recovery of \(\psi\) from \(R_\psi\) is a separate phase retrieval question.

\noindent\begin{minipage}{\linewidth}
\textbf{Theorem 9.4 (sign expansion).} For the bump family of Section 1 set
\(U_i=2M\varepsilon\int_x^y O_i(z)\,dz\) and \(\sigma_S=\prod_{i\in S}\sigma_i\). Then

\[J_m(v_\sigma)=\frac{1}{2M^2}\iint_{x<y}e^{-2M\int_x^y e}\cosh\!\left(\sum_i\sigma_iU_i\right)\,dx\,dy=\sum_{|S|\ {\mathrm{even}}}\sigma_S c_S(m),\]

\[c_S(m)=\frac{1}{2M^2}\iint_{x<y}e^{-2M\int_x^y e}\prod_{i\in S}\sinh U_i\prod_{i\notin S}\cosh U_i\,dx\,dy.\]
\end{minipage}\par

Reflection of \([x,y]\) replaces the exponential of the signed sum by its hyperbolic cosine
after averaging. Expanding \(\cosh(\sum_i\sigma_i U_i)\) gives the formula. The numbers
\(U_i\) all have the same weak sign on each interval because the primitives of the \(O_i\)
are even and nondecreasing in \(|x|\). In particular every coefficient with two indices is
positive.

For a nonempty even set \(S\), let \(\ell_S\) be the round Mercator distance between the
facing edges of its innermost and outermost supports. Then
\(c_S(m)=O(M^{-2}e^{-2M\ell_S})\). In fact, an interval contributing to \(c_S\) crosses a
Mercator distance at least \(\ell_S\), and \(e-\varepsilon\sum_i|O_i|=v_0\) gives the bound
after estimating the hyperbolic functions by exponentials.

For consecutive bumps \(i,i+1\) the logarithmic decay rate of \(c_{\{i,i+1\}}(m)\) is
exactly their gap \(g_i\). For the lower bound, take \(x,y\) just inside the facing support
edges on the northern half. Over a sufficiently small such rectangle no other support meets
\([x,y]\), both active odd integrals have positive lower bounds there, and
\(\int_x^y v_0<g_i+\eta\). Removing the growing exponentials from the two sines leaves the
density \(v_0\) on these intervals and two factors tending uniformly to one. Thus the
integral over that rectangle is bounded below by a positive constant times
\(M^{-2}e^{-2M(g_i+\eta)}\). This proves the rate for every \(\eta>0\).

The corresponding autocorrelations satisfy
\(R_{\psi_\sigma}=\sum_{|S|\ {\mathrm{even}}}\sigma_S\rho_S\), where
\(\rho_S=2^{-N}\sum_\sigma\sigma_S R_{\psi_\sigma}\). Consequently

\[4m^2c_S(m)=\int_{\mathbb R}e^{-2M|h|}\rho_S(h)\,dh.\]

For nonempty even \(S\), \(\rho_S\) vanishes on \(|h|<\ell_S\). To see this, test the lag
measures against an even continuous function supported in this interval. For each \(x<y\),
the sign average cancels whenever \(\int_x^y O_i=0\) for some \(i\in S\). Otherwise the
endpoints span the extreme supports of \(S\) and
\(\int_x^y v_\sigma\ge\int_x^y v_0\ge\ell_S\) for every \(\sigma\). In both cases the
resulting contribution is zero.

\textbf{Theorem 9.5 (distinct gaps).} For fixed bump shapes and amplitudes, pairwise
distinct consecutive Mercator gaps suffice for the high-weight sectors to determine the
hemisphere arrangement up to global reflection.

\textbf{Proof.} Compare \(\sigma\) and \(\gamma\) with \(\gamma\ne\pm\sigma\). Call a gap a
boundary gap if the adjacent entries of \(\sigma_i\gamma_i\) have opposite signs. If an even
set \(S\) has \(\sigma_S\ne\gamma_S\), it contains an odd number of flipped and of unflipped
signs. Its extreme indices therefore span a boundary gap. Let \(g_i\) be the smallest
boundary gap. Under the stated distinctness assumption, only \(S=\{i,i+1\}\) among the
sign-changing sets can have \(\ell_S=g_i\). Every other such set spans a strictly greater
distance. The bounds and the matching logarithmic rate in Theorem 9.4 show that the ratio of
\(J_m(v_\sigma)-J_m(v_\gamma)\) to \(2\sigma_i\sigma_{i+1}c_{\{i,i+1\}}(m)\) tends to one.
Thus the two arrangements differ for every sufficiently large \(M\).

The same proof applies whenever the smallest boundary gap for the chosen comparison is
unique. In the sandwich \(B<A<D\) of Section 5, with the factor \(r\) included in the shape
of the outer bump, the difference is \(2(c_{\{B,A\}}(m)-c_{\{A,D\}}(m))\). Equal gap lengths
can leave competing terms of the same exponential order.

\section{10. Two gap lengths determine the three-bump
arrangement}\label{two-gap-lengths-determine-the-three-bump-arrangement}

For two candidate sign patterns, a boundary gap separates a bump whose sign changes from a
consecutive bump whose sign stays fixed. Theorem 9.5 uses Mercator lengths to distinguish
patterns with a unique shortest boundary gap. Theorem 8.2 supplies a second way to compare
two gaps.

\textbf{Lemma 10.1 (Mercator distortion).} Let \(v>0\) be strictly increasing on \((0,1)\).
For a nonempty interval \(I\), write \(\mathrm{Lat}(I)=\int_I\sqrt v\,dx\) and
\(\mathrm{Merc}(I)=\int_I v\,dx\). If \(I_1\) lies strictly to the left of \(I_2\) with
separated closures, then

\[\frac{\mathrm{Merc}(I_1)}{\mathrm{Lat}(I_1)}<\frac{\mathrm{Merc}(I_2)}{\mathrm{Lat}(I_2)}.\]

\textbf{Proof.} Each ratio is the average of \(\sqrt v\) on its interval with respect to the
positive weight \(\sqrt v\,dx\). Every value on the left interval is smaller than every
value on the right. Equal latitude lengths thus imply
\(\mathrm{Merc}(I_1)<\mathrm{Merc}(I_2)\), and equal Mercator lengths imply
\(\mathrm{Lat}(I_1)>\mathrm{Lat}(I_2)\).

\textbf{Theorem 10.2 (the middle bump).} Let \(B<A<D\) be fixed nonzero, nonnegative bumps
with separated supports on a round background. Keep \(B\) in the north and \(rD\) in the
south for \(r>0\), and switch \(A\) from north to south, with positive amplitudes giving
admissible metrics. The resulting \(v_N,v_S\) have different equivariant spectra.

\textbf{Proof.} Let \(I_1\) be the open gap between \(B\) and \(A\), and \(I_2\) the gap
between \(A\) and \(D\). Use the actual support edges. Each interpolating profile equals
\(v_0\) on both gaps. Unequal latitude lengths separate the rotation-invariant spectra by
Theorem 8.2. For equal latitude lengths, Lemma 10.1 gives
\(\mathrm{Merc}(I_1)<\mathrm{Merc}(I_2)\). Theorem 9.4 gives

\[J_m(v_N)-J_m(v_S)=2\big(c_{BA}(m)-c_{AD}(m)\big).\]

The two coefficients are positive, with logarithmic decay rates given by their respective
Mercator gaps. Thus their difference is strictly positive for all sufficiently large
\(|m|\).

\textbf{Corollary 10.3 (all placements of three bumps).} For three fixed nonzero,
nonnegative bumps of fixed positive amplitudes on a round background, ordered by their
separated northern support intervals, the equivariant spectrum determines their hemisphere
placements up to simultaneous reflection.

\textbf{Proof.} Up to reflecting one configuration globally, any two inequivalent placements
differ at just one bump. Call that bump \(A\) and the others \(B,D\). If the two unchanged
bumps lie in the same hemisphere, the second sector trace changes by
\(\pm2(c_{BA}+c_{AD})\), which is nonzero for every \(m\ne0\). If they lie in opposite
hemispheres and \(A\) is the middle bump, Theorem 10.2 applies. If \(A\) is an end bump, its
coupling to its adjacent bump has logarithmic rate equal to their Mercator gap. Its coupling
to the farther bump decays at least at the rate associated with the full span between their
facing edges, which is strictly longer. The adjacent coupling therefore dominates for large
\(|m|\) by Theorem 9.4.

\textbf{Proposition 10.4 (comparison of gap lengths).} For positive continuous functions
\(f,g\) and separated ordered gaps \(I_1<I_2\), strict monotonicity of \(f/g\) implies that
the two gaps cannot have equal lengths for both densities \(f\) and \(g\). Indeed,
\((\int_I f)/(\int_I g)\) is a weighted average of \(f/g\). For strictly increasing \(v\),
it follows that \(\int_{I_1}v^p\,dx=\int_{I_2}v^p\,dx\) for at most one real exponent \(p\).
For \(c\ge0\) the same conclusion holds for at most one member of the family
\(f_c(v)=\sqrt{v^2+c^2v}\), since \(f_c/f_{c'}=\sqrt{(v+c^2)/(v+c'^2)}\) is strictly
monotone in \(v\) whenever \(c\ne c'\).

For a fixed finite bump family, the two gap estimates also distinguish configurations with
at most two boundary gaps, provided at least one sign is fixed and at least one changes. For
the latitude estimate, use the odd part of the fixed bumps as \(P\) and that of the flipped
bumps as \(Q\) in Lemma 8.1. Any nonadjacent cross-pair spans a boundary gap and has
strictly greater separation than that gap. Consequently, a unique shortest latitude boundary
gap gives the same conclusion as in Theorem 8.2. The Mercator assertion follows from the
proof of Theorem 9.5. If two boundary gaps tie in one coordinate, Lemma 10.1 separates them
in the other.

With three ordered boundary gaps \(I_0<I_1<I_2\), Lemma 10.1 permits the two ties

\[\mathrm{Lat}(I_1)=\mathrm{Lat}(I_2)<\mathrm{Lat}(I_0),\]

\[\mathrm{Merc}(I_0)=\mathrm{Merc}(I_1)<\mathrm{Merc}(I_2).\]

Four bumps with the first and third flipped already have three boundary gaps. Proposition
10.4 gives a length comparison for the further spectral test described in Section 11.

\section{11. Results and questions}\label{results-and-questions}

The table records which spectral data distinguish the profiles considered here. In the
three-bump results the shapes, absolute amplitudes and locations on the northern half are
fixed. A simultaneous reflection always preserves the spectrum.

{% Unnumbered results table.
\begin{longtable}[]{@{}
  >{\raggedright\arraybackslash}p{(\linewidth - 2\tabcolsep) * \real{0.3900}}
  >{\raggedright\arraybackslash}p{(\linewidth - 2\tabcolsep) * \real{0.6100}}@{}}
\toprule\noalign{}
\begin{minipage}[b]{\linewidth}\raggedright
Result
\end{minipage} & \begin{minipage}[b]{\linewidth}\raggedright
Conclusion
\end{minipage} \\
\midrule\noalign{}
\endhead
\bottomrule\noalign{}
\endlastfoot
Local semiclassical invariants (2.1) & All coefficients agree for the reflected two-bump
pair. \\
Inverse traces and height coordinates (3.1--4.1) & The positive quadratic form separates two
bumps. The cubic outer weights give the coupling between distinct levels. \\
Third-trace comparison (5.1--5.3) & The two exterior orderings have strict signs. In the
sandwich construction at least one smooth balancing position exists. \\
Finite data (6.1--7.1) & Odd multiplicity recovers the invariant sector. Any prescribed
finite prefix can tie in a sufficiently large family. Each fixed family has a finite
cutoff. \\
Invariant-sector separation (8.1--8.3) & The mixed derivative and the uniform upper and
lower Green kernel bounds give the exact logarithmic gap rate. A balanced triple is
distinguished by a later trace. \\
Nonzero sectors (9.1--9.3) & A single nonzero weight distinguishes the two-bump pair. The
facing support edges give the logarithmic decay rate. \\
Several bumps (9.4--9.5) & The sign coefficients obey a gap bound. A unique shortest
boundary gap separates the two arrangements at all sufficiently large weights. \\
Three bumps (10.1--10.3) & The equivariant spectrum determines the arrangement up to global
reflection, including equal-gap cases. \\
\end{longtable}
}

\textbf{Four or more bumps.} The first case outside the two-gap argument has four bumps with
the first and third flipped. One possible approach uses the rays \(\kappa=c|m|\), \(c>0\),
with length density \(\sqrt{v^2+c^2v}\). Proposition 10.4 implies that the lengths of
finitely many disjoint gaps are pairwise distinct for all but finitely many choices of
\(c\). To turn this observation into spectral separation one would need an observable with a
nonvanishing leading contribution from the shortest boundary gap and a uniform decay
estimate along that ray.

\textbf{Exact invariant-sector isospectrality.} Can two profiles in a common \((W,|a'|)\)
class have the same full \(m=0\) spectrum without being related by reflection? A uniform
cutoff \(K_*\le n\) in a family to which Theorems 6.2 and 7.1 both apply would give such a
pair. Existence of one exactly isospectral pair would supply no such bound for the whole
family.

\textbf{Accumulating swaps and signed perturbations.} If infinitely many gaps accumulate,
there may be no shortest gap. A limiting argument for the finite-trace constructions would
also have to keep the two limiting profiles distinct up to reflection. Allowing bumps to
change sign raises a separate question, since the weights in the lower bounds of Sections 8
and 9 depend on positivity.

\textbf{Ordinary spectrum.} Does the unlabeled Laplace spectrum distinguish the sandwich
arrangements when their two latitude gaps tie? The proof in Section 10 uses nonzero weight
labels in that case.

\textbf{Smooth single wells.} Does the asymptotic equivariant spectrum of [1, Definition
2.1] determine every smooth single well? Section 2 isolates the branch-labeling issue in
the proof. Section 9 resolves the exact large-weight comparison for the explicit two-bump
pair. For a general profile one can also ask which swaps preserve the autocorrelation of
\(1/v\) in the Mercator coordinate, the information carried by all the \(J_m\).

\textbf{Relation to cyclic covers.} For a rotation acting through an irrational angle on a
surface of revolution, the equivariant spectrum is the cross-section data obtained from the
full cyclic tower by recovering the twisted heat traces and Fourier decomposing the rotation
action. The results above describe what these data distinguish within the fixed bump
families. Recovering an arbitrary smooth profile requires a more general uniqueness
argument.

\section{References}\label{references}

[1] E. B. Dryden, D. Macedo and R. Sena-Dias. Recovering \(S^1\)-invariant metrics on
\(S^2\) from the equivariant spectrum. \emph{International Mathematics Research Notices}
2016, no. 16, 4882--4902. \href{https://arxiv.org/abs/1501.02830v2}{arXiv:1501.02830v2}.

[2] Y. Colin de Verdière. A semi-classical inverse problem II: reconstruction of the
potential. 2008. \href{https://arxiv.org/abs/0802.1643}{arXiv:0802.1643}.

[3] M. Engman. The spectrum and isometric embeddings of surfaces of revolution.
\emph{Canadian Mathematical Bulletin} 49 (2006), no. 2, 226--236.
\href{https://doi.org/10.4153/CMB-2006-023-7}{doi:10.4153/CMB-2006-023-7}.

[4] J. Brüning and E. Heintze. Spektrale Starrheit gewisser Drehflächen.
\emph{Mathematische Annalen} 269 (1984), 95--101.
\href{https://doi.org/10.1007/BF01455998}{doi:10.1007/BF01455998}.

\end{document}
